\titledquestion{Arrays}

\textbf{Arrays}

We discussed \textit{sequences} in lecture, which are immutable arrays that
admit nice parallel properties due to having constant-time access to each
index.

Sequences, however, have some memory disadvantages. For instance, to map a function
on a sequence, we need to allocate a brand new sequence. For some applications,
this can be rather disadvantageous.

Imperative arrays can perform \textit{in-place} operations, which means performing
operations on an array \textit{without} having to allocate a brand new array. In
other words, by reusing the same space in place, via mutation.

Sequences are not naturally in-place. But what if we combine them with reference
cells? Consider the following type:

\begin{codeblock}
  type 'a array = 'a ref Seq.t
\end{codeblock}

Now we have recovered arrays from sequences.

\begin{parts}
  \part[5]\hfill

    Implement the following function:
    \spec
      {update}
      {\code{'a array -> int * 'a -> unit}}
      {\code{true}}
      {\code{update A (i, x)} updates the \code{i}th index of the array \code{A} to
      the element \code{x}. Raises \code{Range} if out of range.}

    Assume that there exists an exception \code{Range} declared somewhere.

    \begin{solutionorbox}[15em]
      \begin{codeblock}
        fun update A (n, x) =
          if n >= Seq.length x orelse n < 0 then
            raise Range
          else
            Seq.nth A n := x
      \end{codeblock}
    \end{solutionorbox}

  \newpage

  \part[5]\hfill

    Implement the following function:
    \spec
      {tabulate}
      {\code{(int -> 'a) -> int -> 'a array}}
      {\code{f} is total, \code{n >= 0}}
      {\code{tabulate f n} evaluates to a fresh array of length \code{n},
      where the \code{i}th index is equivalent to \code{f i}}

    \begin{solutionorbox}[15em]
      \begin{codeblock}
        fun tabulate f n =
          Seq.tabulate (fn i => ref (f i)) n
      \end{codeblock}
    \end{solutionorbox}

  \part[5]\hfill

    Implement the following function:
    \spec
      {map}
      {\code{('a -> 'a) -> 'a array -> unit}}
      {\code{f} is total}
      {\code{map f A} maps all the elements \code{x} in \code{A} to \code{f x},
      in-place}

    You may assume the existence of
    \code{Seq.iter : ('a -> unit) -> 'a Seq.t -> unit}, which simply
    applies a function to each element of a sequence.

    \begin{solutionorbox}[15em]
      \begin{codeblock}
        fun map f A =
          (Seq.map (fn r => r := f (!r)) A; ())
      \end{codeblock}
    \end{solutionorbox}

  \newpage

  \part[12]\hfill

    Now, let's do a fun use case for in-place algorithms: in-place sorting.\footnote{
      Any reasonable $O(n \log n)$ algorithm would be absurd to ask you to implement
      during an exam, especially in-place. But we can do a basic one. Let's do bubblesort.
    }

    Bubblesort is a sorting algorithm which is done by repeatedly swapping
    out-of-order adjacent elements, until the array is completely sorted. We can
    describe the algorithm as follows:

    \begin{lstlisting}[style=15150code, numbers=none]
      for adjacent pairs in array:
        if pair is out of order:
          swap them

      if any pairs were swapped, go back to loop
    \end{lstlisting}

    Implement the following function:
    \spec
      {bubsort}
      {\code{('a * 'a -> order) -> 'a array -> unit}}
      {\code{cmp} is a valid comparison function}
      {\code{bubsort cmp A} sorts the array \code{A} via \code{cmp}, in-place}

    You may use the following \code{swap : 'a ref -> 'a ref -> unit} function:

    \begin{codeblock}
      fun swap r1 r2 =
        let
          val r1_v = !r1
        in
          ( r1 := !r2; r2 := r1_v )
        end
    \end{codeblock}

    We give you the following skeleton of code, where
    \code{adj_pairs : 'a array -> ('a ref * 'a ref) list} obtains all of the
    adjacent pairs within a list. So for instance,
    \code{adj_pairs} $\langle \code{r1}, \code{r2}, \code{r3} \rangle$ $\eeq$
    \code{[(r1, r2), (r2, r3)]}.

    \begin{codeblock}
      fun bubsort cmp A =
        let
          val pairs = adj_pairs A

          fun aux pairs = (* Implement me! *)
        in
          aux pairs
        end
    \end{codeblock}

    \newpage

    \begin{solutionorbox}[45em]
      \begin{codeblock}
        fun aux pairs =
          if
            List.foldl (fn ((r1, r2), acc) =>
              case cmp (!r1, !r2) of
                GREATER => (swap r1 r2; true)
              | _ => acc
            ) false pairs
          then
            aux pairs
          else
            ()
      \end{codeblock}

    \end{solutionorbox}
\end{parts}

\newpage
